\documentclass[11pt]{article}
\usepackage[a4paper,margin=26mm]{geometry}
\usepackage{amsmath,amssymb,amsthm,booktabs}
\setlength{\emergencystretch}{2em}
\newtheorem{theorem}{Theorem}
\title{Every rank in the alternating group interval}
\author{Human Compression project: principal recovery of L038}
\date{Author integration; distinct review pending}
\begin{document}\maketitle

\begin{theorem}
For every integer $n\ge12$ and every integer
$3\le r\le\lfloor(n-1)/2\rfloor$, the alternating group $A_n$
has a string C-group representation of rank $r$.
\end{theorem}

This is Theorem~1.1 of Fernandes and Leemans,
\emph{String C-group representations of alternating groups},
in the fixed edition arXiv:1810.12450v2, 4 January 2019. The two
companion proofs establish the required permutation constructions.
Here we assemble their conclusions with the already acquired
contraction theorem and two existing degree-twelve witnesses.
The maximality of the interval is a separate earlier theorem;
it is not asserted as a new conclusion of this recovery.

\section{The precise acquired interfaces}
A string C-marking $(s_0,\ldots,s_{R-1})$ is irreducible if every
adjacent product has order greater than two. Put
$p_i=|s_{i-1}s_i|$. The acquired proof of Brooksbank--Leemans
rank contraction gives the following statement: if
$p_3,\ldots,p_{t+2}$ are odd, where $0\le t\le R-3$, the same
group has C-markings of all ranks $R-t,\ldots,R$. The ordering
may first be reversed. Every contraction preserves irreducibility.
We use this proved indexing, rather than the extraneous $p_2$
test in the source's printed iteration formula.

Two earlier maximal-rank constructions are retained as inputs.
They are exactly the inputs already recorded in the acquired
L037 application \emph{Rank intervals from one string marking},
Section~2, and are explicitly attributed to the earlier literature
in the fixed L038 source:
\begin{itemize}
\item For $n\equiv1\pmod4$, $n\ge13$, rank $R=(n-1)/2$,
with type $[10,3^{R-2}]$. The source recalls the prior alternating
construction in Theorem~3.2 and uses it in Theorem~4.6.
\item For even $n\ge14$, rank $R=(n-2)/2$, with type
$[5,6,3^{R-6},6,6,3]$. This is the earlier maximal-rank
construction explicitly recalled at the start of the proof of
source Theorem~4.1.
\end{itemize}
These include generation and the full intersection property;
they are not merely assertions about adjacent orders. Contraction
therefore supplies every rank $3,\ldots,R$ in the first case
and every rank $6,\ldots,R$ in the second. The latter uses
exactly $R-6$ consecutive odd bonds beginning at $p_3$; when
$R=6$, no contraction is required.

Our reviewed companion \emph{Finite end groups for alternating
group constructions} proves, with faithful finite end certificates,
the following four families from the source:
\begin{center}\small
\begin{tabular}{llll}
\toprule
Degree & Rank & Type & Source theorem\\\midrule
$n\equiv2\ (4),\ n\ge10$ &5&$[5,5,6,n-5]$&4.3\\
$n\equiv0\ (4),\ n\ge12$ &5&$[3,4,6,n-7]$&4.5\\
$n\equiv0\ (4),\ n\ge16$ &4&$[3,12,\operatorname{lcm}(n-8,6)]$&4.4\\
$n\equiv3\ (4),\ n\ge15$ &6&$[n-10,6,3,5,3]$&4.13\\\bottomrule
\end{tabular}
\end{center}
In the first row reverse the marking: its tested bonds are five
and five, giving ranks four and three. In the third row reversal
gives third bond three and hence rank three. In the fourth row
the tested bonds are three, five and three, giving ranks five,
four and three. All adjacent orders are greater than two, so the
irreducibility hypothesis is present in each application.

The second companion, \emph{Signed permutation groups at the ends
of alternating group constructions}, proves the exact source
Theorems~4.8--4.10 by two uniform signed-group calculations and
explicit parity intersections. Its conclusions are:
\begin{itemize}
\item If $r\ge7$ is odd and $k\ge1$, $A_{2r+1+4k}$ has a
rank-$r$ C-marking. The case $k=0$ is the earlier maximal-rank
seed, explicitly retained through source Theorem~4.7.
\item If $r\ge8$ is even and $k\ge0$, $A_{2r+3+4k}$ has a
rank-$r$ C-marking, including the $k=0$ boundary.
\end{itemize}
The companion proves their exact displayed types as well.
It uses the earlier odd maximal-rank seed only where it explicitly
appears, and does not assume the principal theorem being recovered.

\section{Degree twelve from the existing catalogue}
Two tuples from the already accepted L071 alternating-group census
are enough for the remaining degree-twelve boundary. They are
listed here so the existence assertion is independently readable:
\begin{align*}
 a_0&=(5\ 6)(7\ 8)(9\ 10)(11\ 12),\\
 a_1&=(3\ 5)(4\ 6)(8\ 9)(10\ 11),\\
 a_2&=(1\ 3)(2\ 4)(5\ 7)(6\ 8),
\end{align*}
and
\begin{align*}
 b_0&=(9\ 10)(11\ 12),\\
 b_1&=(5\ 6)(7\ 9)(8\ 11)(10\ 12),\\
 b_2&=(1\ 5)(2\ 6)(3\ 7)(4\ 8)(9\ 10)(11\ 12),\\
 b_3&=(2\ 3)(7\ 8)(9\ 11)(10\ 12).
\end{align*}
The recorded ranks are three and four, with types $[12,20]$ and
$[6,8,6]$. Every mark is even, each tuple generates a group of
order $12!/2$, and all $8^2=64$ and $16^2=256$ ordered pairs of
subset subgroups satisfy the intersection property. Their string
commutations also hold. Thus these are representations of $A_{12}$.
The accompanying small certificate checks these two already
recorded objects, without performing a new classification. In
the retained L071 records they are rank-three row 43, involution
indices $(117,1903,23898)$, and rank-four row 20, indices
$(1,122,32811,5368)$. The exact predecessor record file and its
accepted review are bound in the dependency manifest. Rank five
at $n=12$ is already the second low-rank family above.

\section{Completion of the interval}
\begin{proof}[Proof of Theorem 1]
For $n=12$, the two recorded witnesses and the rank-five family
give all three required ranks. If $n\equiv1\pmod4$, the first
maximal-rank seed and contraction give the whole interval.

Suppose next that $n\ge14$ is even. The second maximal-rank
seed and contraction give all ranks at least six through
$R=(n-2)/2$. When $n\equiv2\pmod4$, the first low-rank
family and its two contractions give ranks three, four and five.
When $n\equiv0\pmod4$, necessarily $n\ge16$; the second
and third low-rank families give rank five and ranks four and
three, respectively. These intervals meet and cover every
$3\le r\le R$.

Finally let $n\equiv3\pmod4$, so $n\ge15$ and
$R=(n-1)/2$ is odd. The rank-six family and its contractions
give ranks three through six. For every odd $7\le r\le R$,
put $k=(R-r)/2$. This is a nonnegative integer and
$n=2r+1+4k$; use the higher odd family if $k\ge1$ and the
retained maximal-rank seed if $k=0$. For every even
$8\le r\le R-1$, put $k=(R-1-r)/2$. Again $k$ is a
nonnegative integer, and $n=2r+3+4k$; the higher even family
applies. If $R=7$, the even list is empty, as required.
Together these give every rank from three through $R$.
\end{proof}

\section{What is recovered and what is retained}
The whole principal existence statement, source Theorem~1.1,
is recovered. The constructions of source Theorems~4.3--4.5,
4.8--4.10 and 4.13 are recovered with their exact types by the
two companions. For source Theorems~4.1 and 4.6, this integration
gives their rank-existence conclusions from earlier maximal seeds.
For 4.2, 4.11, 4.12 and 4.14, it likewise gives the required
rank existence, using contraction instead of the source's
different lower-rank diagrams. Their exact types and every
auxiliary proposition are not claimed as additional recoveries.

The earlier maximal constructions, standard end criterion and
polytope correspondence retain their original attribution.
The rank-contraction theorem is the acquired proof of
Brooksbank--Leemans, not a new theorem in this integration.
Our work lies in the complete end calculations, the faithful
finite reductions with their constraints, and the assembly
which eliminates separate lower-rank constructions. This is
not an open-conjecture solution or a historical priority claim.
No new Lean formalization or complexity measurement is part
of this intake recovery.

\begin{thebibliography}{9}
\bibitem{FL} M. E. Fernandes and D. Leemans,
\emph{String C-group representations of alternating groups},
arXiv:1810.12450v2, 4 January 2019; Ars Math. Contemp. 17 (2019),
291--310.
\bibitem{BL} P. A. Brooksbank and D. Leemans,
\emph{Rank reduction of string C-group representations},
arXiv:1812.01055v2, 3 May 2019; Proc. Amer. Math. Soc. 147
(2019), 5421--5426.
\bibitem{LM} D. Leemans and M. Mixer,
\emph{Algorithms for classifying regular polytopes with a fixed
automorphism group}, Contributions to Discrete Mathematics
7(2) (2012), 105--118. The acquired L071 census supplies the
two literal witnesses above.
\end{thebibliography}
\end{document}
